\documentclass[11pt,letterpaper]{article}
% KJO_preamble.tex
% Shared LaTeX preamble for all KJO documentation documents.
% \input this file immediately after \documentclass{}.

% --- Encoding and fonts -------------------------------------------------------
\usepackage[utf8]{inputenc}
\usepackage[T1]{fontenc}
\usepackage{lmodern}
\usepackage{microtype}

% --- Page geometry ------------------------------------------------------------
\usepackage[
  letterpaper,
  top=1in, bottom=1in,
  left=1.25in, right=1in
]{geometry}

% --- Core utilities -----------------------------------------------------------
\usepackage{amsmath}
\usepackage{graphicx}
\usepackage{xcolor}
\usepackage{parskip}        % space between paragraphs; no indent
\usepackage{enumitem}
\usepackage{caption}
\usepackage{subcaption}
\usepackage{float}

% --- Tables -------------------------------------------------------------------
\usepackage{booktabs}
\usepackage{longtable}
\usepackage{array}
\usepackage{multirow}

% --- Code listings ------------------------------------------------------------
\usepackage{listings}
\lstset{
  basicstyle=\ttfamily\small,
  keywordstyle=\bfseries\color{KJOblue},
  commentstyle=\itshape\color{gray},
  stringstyle=\color{KJOgreen!70!black},
  breaklines=true,
  frame=single,
  framesep=4pt,
  xleftmargin=6pt,
  xrightmargin=4pt,
  numbers=left,
  numberstyle=\tiny\color{gray},
  numbersep=8pt,
  language=C++,
}

% --- TikZ / PGF ---------------------------------------------------------------
\usepackage{tikz}
\usetikzlibrary{
  arrows.meta,
  shapes.geometric,
  shapes.misc,
  shapes.symbols,
  positioning,
  fit,
  calc,
  backgrounds,
  decorations.pathreplacing,
  decorations.markings,
  matrix,
}
\usepackage{pgfplots}
\pgfplotsset{compat=1.18}

% --- Headers / footers --------------------------------------------------------
\usepackage{fancyhdr}
\usepackage{lastpage}
\pagestyle{fancy}
\fancyhf{}
\renewcommand{\headrulewidth}{0.4pt}
\renewcommand{\footrulewidth}{0.4pt}
% Per-document: \lhead{}, \rhead{}, \cfoot{} are set in the main .tex file.

% --- Section formatting -------------------------------------------------------
\usepackage{titlesec}
\titleformat{\section}{\large\bfseries\color{KJOblue}}{{\thesection}}{0.8em}{}[\titlerule]
\titleformat{\subsection}{\normalsize\bfseries\color{KJOblue!80!black}}{{\thesubsection}}{0.8em}{}
\titleformat{\subsubsection}{\normalsize\bfseries}{{\thesubsubsection}}{0.8em}{}

% --- Hyperlinks ---------------------------------------------------------------
\usepackage[
  colorlinks=true,
  linkcolor=KJOblue,
  urlcolor=KJOblue!70!black,
  citecolor=KJOblue,
  pdfborder={0 0 0},
]{hyperref}

% --- Color palette -----------------------------------------------------------
\definecolor{KJOblue}   {RGB}{ 30,  90, 160}   % RCU
\definecolor{KJOred}    {RGB}{180,  30,  30}   % EMU
\definecolor{KJOgreen}  {RGB}{ 20, 120,  60}   % GSEMU
\definecolor{KJOorange} {RGB}{200, 110,  20}   % GSE / pressurization
\definecolor{KJOgray}   {RGB}{100, 100, 100}   % Buses / links
\definecolor{KJOpurple} {RGB}{110,  40, 140}   % Shared libraries
\definecolor{KJObg}     {RGB}{248, 248, 252}   % Light background for code

% --- Convenience macros -------------------------------------------------------

% Hardware component name (small caps)
\newcommand{\hw}[1]{\textsc{#1}}

% Software module name (monospace)
\newcommand{\sw}[1]{\texttt{#1}}

% Command name (monospace bold)
\newcommand{\cmd}[1]{\texttt{\textbf{#1}}}

% Pin / signal name
\newcommand{\sig}[1]{\texttt{#1}}

% Note / callout box
\usepackage{tcolorbox}
\tcbuselibrary{skins}
\newtcolorbox{notebox}{
  enhanced,
  colback=KJOblue!5,
  colframe=KJOblue!40,
  fonttitle=\bfseries,
  title=Note,
  left=4pt, right=4pt, top=2pt, bottom=2pt,
}
\newtcolorbox{warningbox}{
  enhanced,
  colback=KJOred!5,
  colframe=KJOred!40,
  fonttitle=\bfseries\color{KJOred},
  title=Warning,
  left=4pt, right=4pt, top=2pt, bottom=2pt,
}

% Title page macro: \KJOtitlepage{title}{subtitle}{version}{date}{description}
\newcommand{\KJOtitlepage}[5]{%
  \begin{titlepage}
    \centering
    \vspace*{2cm}
    {\Huge\bfseries\color{KJOblue} #1 \par}
    \vspace{0.5cm}
    {\Large\color{KJOgray} #2 \par}
    \vspace{1.5cm}
    \rule{0.5\textwidth}{0.4pt}\par
    \vspace{1cm}
    {\large Ken Overton \par}
    \vspace{0.3cm}
    {\normalsize \textit{KJO Liquid Rocket Motor Control System} \par}
    \vspace{1.5cm}
    \begin{tabular}{ll}
      \textbf{Version:} & #3 \\[4pt]
      \textbf{Date:}    & #4 \\
    \end{tabular}
    \vspace{1.5cm}\par
    \begin{minipage}{0.75\textwidth}
      \centering
      \normalsize #5
    \end{minipage}
    \vfill
    {\small\color{KJOgray} \textcopyright\ 2025--2026 Ken Overton.
     All rights reserved. \par}
  \end{titlepage}%
}

\newcommand{\docversion}{0.1}
\newcommand{\docdate}{September 2026}
\lhead{\textsc{KJO EMU Servo Tuner}}
\rhead{\textsc{User Guide} --- v\docversion}
\cfoot{\thepage\ of \pageref{LastPage}}
\hypersetup{pdftitle={KJO EMU Servo Tuner User Guide}, pdfauthor={Ken Overton}}

\begin{document}
% KJO_tikz_styles.tex
% Shared TikZ node and edge styles for all KJO block diagrams.
% \input this file inside the document body before any tikzpicture environments.

\tikzset{
  % ── Hardware block (peripheral chip / PCB component) ──────────────────────
  hwblock/.style={
    rectangle, draw=KJOgray, line width=1pt,
    fill=KJOgray!10, rounded corners=4pt,
    minimum width=2.6cm, minimum height=0.8cm,
    font=\small, align=center, inner sep=4pt,
  },
  % Sub-variant: MCU (larger, blue border)
  mcublock/.style={
    hwblock,
    draw=KJOblue, fill=KJOblue!8,
    minimum width=3.0cm, minimum height=1.0cm,
    font=\small\bfseries,
  },
  % Sub-variant: power / battery
  powerblock/.style={
    hwblock,
    draw=KJOorange, fill=KJOorange!8,
  },
  % ── Software module box ────────────────────────────────────────────────────
  swmodule/.style={
    rectangle, draw=KJOblue, line width=0.8pt,
    fill=KJOblue!6, rounded corners=3pt,
    minimum width=2.8cm, minimum height=0.7cm,
    font=\small, align=center, inner sep=3pt,
  },
  % Sub-variant: shared library
  sharedlib/.style={
    swmodule,
    draw=KJOpurple, fill=KJOpurple!6,
    dashed,
  },
  % ── Bus / link annotation ──────────────────────────────────────────────────
  busline/.style={
    draw=KJOgray, line width=1.5pt,
    -{Stealth[length=6pt,width=4pt]},
  },
  bidbusline/.style={
    draw=KJOgray, line width=1.5pt,
    {Stealth[length=6pt,width=4pt]}-{Stealth[length=6pt,width=4pt]},
  },
  radioline/.style={
    draw=KJOblue!60, line width=1.5pt, dashed,
    {Stealth[length=6pt,width=4pt]}-{Stealth[length=6pt,width=4pt]},
  },
  gpioline/.style={
    draw=KJOgreen!60!black, line width=1.2pt,
    -{Stealth[length=5pt,width=4pt]},
  },
  % ── Flowchart shapes ──────────────────────────────────────────────────────
  flowproc/.style={
    rectangle, draw=KJOblue!70, line width=0.8pt,
    fill=KJOblue!8, rounded corners=3pt,
    minimum width=3.2cm, minimum height=0.7cm,
    font=\small, align=center, inner sep=3pt,
  },
  flowdec/.style={
    diamond, draw=KJOorange, line width=0.8pt,
    fill=KJOorange!8,
    minimum width=2.2cm, minimum height=0.9cm,
    font=\small, align=center, inner sep=1pt,
  },
  flowterminal/.style={
    rounded rectangle, draw=KJOblue, line width=1pt,
    fill=KJOblue!12,
    minimum width=2.8cm, minimum height=0.7cm,
    font=\small\bfseries, align=center, inner sep=3pt,
  },
  flowarrow/.style={
    draw=KJOgray, line width=0.9pt,
    -{Stealth[length=5pt,width=4pt]},
  },
  % ── Propellant schematic styles ────────────────────────────────────────────
  % Generic pipe line
  pipe/.style={
    draw=black, line width=1.8pt,
  },
  pipearrow/.style={
    pipe,
    -{Stealth[length=7pt,width=5pt]},
  },
  % Oxidizer (N2O) — blue
  oxpipe/.style={
    draw=KJOblue, line width=1.8pt,
  },
  oxpipearrow/.style={
    oxpipe,
    -{Stealth[length=7pt,width=5pt]},
  },
  % Fuel — red
  fuelpipe/.style={
    draw=KJOred, line width=1.8pt,
  },
  fuelpipearrow/.style={
    fuelpipe,
    -{Stealth[length=7pt,width=5pt]},
  },
  % GSE fill line — green
  gsepipe/.style={
    draw=KJOgreen!70!black, line width=1.8pt,
  },
  gsepipearrow/.style={
    gsepipe,
    -{Stealth[length=7pt,width=5pt]},
  },
  % Pressurization — orange dashed
  presspipe/.style={
    draw=KJOorange, line width=1.4pt, dashed,
  },
  presspipearrow/.style={
    presspipe,
    -{Stealth[length=6pt,width=4pt]},
  },
  % Tank (tall rectangle, rounded corners)
  tank/.style={
    rectangle, draw=black, line width=1.5pt,
    rounded corners=4pt, fill=white,
    minimum width=1.8cm, minimum height=2.4cm,
    font=\small\bfseries, align=center, inner sep=4pt,
  },
  n2otank/.style={
    tank, draw=KJOblue, fill=KJOblue!8,
  },
  fueltank/.style={
    tank, draw=KJOred, fill=KJOred!8,
  },
  srctank/.style={
    tank, draw=KJOgreen!70!black, fill=KJOgreen!6,
  },
  % Valve box (servo-actuated)
  valvebox/.style={
    rectangle, draw=black, line width=1pt,
    fill=white, minimum width=2.2cm, minimum height=0.65cm,
    font=\footnotesize\bfseries, align=center, inner sep=3pt,
  },
  gseval/.style={
    valvebox, draw=KJOgreen!70!black, fill=KJOgreen!8,
  },
  emuval/.style={
    valvebox, draw=KJOred, fill=KJOred!6,
  },
  % Component (generic inline component)
  component/.style={
    rectangle, draw=KJOgray, line width=0.8pt,
    fill=white!96!gray, minimum width=2.2cm, minimum height=0.6cm,
    font=\footnotesize, align=center, inner sep=3pt,
  },
  % Instrument callout circle (PT, TT, TC)
  instrument/.style={
    circle, draw=KJOgray, line width=0.8pt,
    fill=white, minimum size=0.6cm,
    font=\tiny\bfseries, align=center, inner sep=0pt,
  },
  % Sensor lead line
  sensorlead/.style={
    draw=KJOgray, line width=0.7pt,
  },
  % Boundary box (dashed outline for GSE / Rocket regions)
  gseboundary/.style={
    draw=KJOgreen!50, line width=1pt, dashed,
    rounded corners=6pt,
  },
  rocketboundary/.style={
    draw=KJOred!40, line width=1pt,
    rounded corners=6pt,
  },
  % Sequence diagram lifeline
  lifeline/.style={
    draw=KJOgray, line width=0.6pt, dashed,
  },
  % Sequence diagram message arrow
  seqarrow/.style={
    draw=KJOblue, line width=0.9pt,
    -{Stealth[length=5pt,width=4pt]},
  },
  % State node (FSM)
  statenode/.style={
    circle, draw=KJOblue, line width=1pt,
    fill=KJOblue!10,
    minimum size=1.2cm,
    font=\small\bfseries, align=center,
  },
  statetransition/.style={
    draw=KJOgray, line width=0.9pt,
    -{Stealth[length=5pt,width=4pt]},
    font=\footnotesize,
  },
}


\KJOtitlepage%
  {EMU Servo Tuner}%
  {Bench Tool User Guide}%
  {\docversion}{\docdate}%
  {A Serial-console tool, flashed onto the actual EMU hardware, for
   interactively tuning the PWM open/close and potentiometer open/close
   limits of each of EMU's ball valves.}

\tableofcontents\newpage

%═══════════════════════════════════════════════════════════════════════════════
\section{Purpose}
%═══════════════════════════════════════════════════════════════════════════════

EMU Servo Tuner is a standalone sketch, built and flashed onto the same
physical EMU hardware (Feather M0 + Servo FeatherWing) as production
firmware, that lets an operator select one of a small built-in list of
valves and interactively tune its PWM/encoder limits over Serial -- the
same target/adjust workflow as the standalone \textbf{Servo Tuner}, but
driven by a per-valve table instead of one hardcoded channel/pin pair.

\begin{warningbox}
The valve table built into this tool (§\ref{sec:valve-table}) is a
snapshot taken when the tool was written -- it is \textbf{not}
automatically kept in sync with EMU's production \texttt{KJO\_Valve.h}.
Treat the printed limits as a starting point only, and always verify
against the current production reference doc (\textit{KJO EMU
Reference}) before trusting a number from this tool's menu. In
particular, this tool's built-in list still includes a \texttt{FILL}
entry; production EMU has no Fill valve of its own (it lives on GSEMU) --
the entry's numbers actually match GSEMU's Fill valve calibration, not an
EMU valve.
\end{warningbox}

%═══════════════════════════════════════════════════════════════════════════════
\section{Hardware Required}
%═══════════════════════════════════════════════════════════════════════════════

\begin{itemize}[nosep]
  \item EMU hardware: Adafruit Feather M0 + PWM/Servo FeatherWing (PCA9685)
  \item The valve(s) to be tuned, wired to their normal EMU pin/channel
        assignments
  \item USB cable for Serial (115200 baud)
\end{itemize}

%═══════════════════════════════════════════════════════════════════════════════
\section{Built-In Valve Table}
\label{sec:valve-table}
%═══════════════════════════════════════════════════════════════════════════════

\begin{center}
\begin{tabularx}{\textwidth}{@{}lllllY@{}}
\toprule
Valve & Servo ch. & PWM open & PWM close & Pot pin & Encoder open/closed \\
\midrule
FILL      & 4 & 2329 & 3918 & A4 & 1310 / 3004 \\
OX FEED   & 5 & 2050 & 3865 & A3 & 1146 / 2846 \\
FUEL FEED & 6 & 2280 & 3950 & A1 & 1217 / 2887 \\
\bottomrule
\end{tabularx}
\end{center}

These are the limits displayed when a valve is selected -- they are the
tool's own starting reference, not necessarily what is currently flashed
to production EMU.

%═══════════════════════════════════════════════════════════════════════════════
\section{How to Use}
%═══════════════════════════════════════════════════════════════════════════════

Open a Serial monitor at 115200 baud after powering the board.

\subsection{Outer Loop -- Valve Selection}

\begin{enumerate}[nosep]
  \item The tool prints a numbered menu of the valves in
        §\ref{sec:valve-table}. Enter the number and press Enter.
  \item It prints that valve's servo channel, pot pin, and current
        PWM/encoder open-close limits.
\end{enumerate}

\subsection{Inner Loop -- Target / Adjust}

Identical workflow to the standalone Servo Tuner, using the selected
valve's own channel and pot pin:

\begin{sloppypar}
\begin{enumerate}[nosep]
  \item \textbf{Target}: enter a PWM count (the valve's own open/close
        values are printed as a range hint).
  \item \textbf{Hysteresis delta}: enter a count value; the tool
        approaches the target from \texttt{target~+~delta} counts away
        to remove backlash from the reading.
  \item The servo moves and the resulting potentiometer reading is
        printed.
  \item \textbf{Action prompt} -- type one letter and press Enter:
        \begin{itemize}[nosep]
          \item \texttt{a} (adjust): enter a signed delta; the tool moves
                away from the \emph{previous} target by the hysteresis
                amount, then to the new target, and prints the new pot
                reading.
          \item \texttt{c} (continue): return to step 1 for a fresh
                target on the \emph{same} valve.
          \item \texttt{f} (finished): return to the outer valve-selection
                menu.
        \end{itemize}
\end{enumerate}
\end{sloppypar}

%═══════════════════════════════════════════════════════════════════════════════
\section{Notes}
%═══════════════════════════════════════════════════════════════════════════════

\begin{itemize}[nosep]
  \item The Serial console has no timeout -- the tool waits indefinitely
        at every prompt.
  \item Backspace is supported while typing a value.
  \item An unrecognised action-prompt letter simply re-prints the prompt;
        it does not exit or lose the current target.
\end{itemize}

%═══════════════════════════════════════════════════════════════════════════════
\section{Version History}
%═══════════════════════════════════════════════════════════════════════════════
\begin{center}
\begin{tabularx}{\textwidth}{@{}llY@{}}
\toprule
Version & Date & Notes \\
\midrule
0.1 & September 2026 & Initial user guide. \\
\bottomrule
\end{tabularx}
\end{center}

\end{document}
